% GENERATED, edit the French unit or the translation table
\mxnewpage\section*{Solutions EV1-B}\label{EV1-B-CORR}
\textbf{MCQ:} 1 B, 2 B, 3 D, 4 A, 5 D, 6 C, 7 C, 8 B, 9 A, 10 C, 11 B, 12 C. Each correct answer: \num{0,5} point; total 6.
\paragraph{Checks of the fluency questions.} Order of operations; sum of fractions; subtraction of exponents; multiplying coefficient \num{1,2}; division by \num{0,8}; product \num{1,25} times \num{0,8}; solving an affine equation; vertical increment divided by horizontal; mean; angle divided by 360; complement to 1; substitution in brackets.
\subsection*{Exercise 1 -- Growth by pattern}
\begin{enumerate}
\item $3,8,15$. \textit{Marking: 1 point. Result and basic justification expected; half a point possible for a correct but unfinished approach.}
\item Difference $2n+3>0$: strictly increasing. \textit{Marking: 2 points. Method or justification: 1; calculation and conclusion: 1.}
\item $n^2+2n-99=(n-9)(n+11)\geq0$ with $n\geq1$: $n\geq9$; $u_9=99$. \textit{Marking: 2 points. Method or justification: 1; calculation and conclusion: 1.}
\end{enumerate}
\subsection*{Exercise 2 -- Two dimensions and a constraint}
\begin{enumerate}
\item $x+y=13$, $xy=40$, x,y positive. \textit{Marking: 1 point. Result and basic justification expected; half a point possible for a correct but unfinished approach.}
\item Roots of $t^2-13t+40=(t-5)(t-8)$: 5 m and 8 m. \textit{Marking: 2 points. Method or justification: 1; calculation and conclusion: 1.}
\item $A(x)=x(13-x)=\dfrac{169}{4}-(x-\dfrac{13}{2})^2$, maximum $\dfrac{169}{4}=\num{42,25}$ m$^2$. Difference $\dfrac{9}{4}$ m$^2$. \textit{Marking: 2 points. Method or justification: 1; calculation and conclusion: 1.}
\end{enumerate}
\subsection*{Exercise 3 -- Angle then approximation}
\begin{enumerate}
\item $\cos \theta =\dfrac{3}{6}=\dfrac{1}{2}$, so $60^\circ$. \textit{Marking: 1 point. Result and basic justification expected; half a point possible for a correct but unfinished approach.}
\item Square $4+9-6=7$; norm $\sqrt7$. \textit{Marking: 1 point. Result and basic justification expected; half a point possible for a correct but unfinished approach.}
\item Approximation \num{1,06}; exact \num{1,0609}; error \num{0,0009}. This independent question checks local differentiation. \textit{Marking: 2 points. Method or justification: 1; calculation and conclusion: 1.}
\end{enumerate}
\paragraph{Calibration.} Target duration: 20--25 min for the fluency questions, 80--85 min for the exercises, 10--15 min for proofreading. These durations are editorial estimates to be measured in class. Accept any correct method; do not penalise twice a numerical error that is propagated consistently. The proposed marking scheme is not a decision of the team to harmonise marking.
